\section{Discussion}

The full-release analysis did not detect the predicted Condition $\times$ instructed Phase $\times$ SOA interaction. Its inference is provisional: two of the 42 participant records contain identical behavioural recordings with reversed intervention labels and different questionnaire responses. Resolving that source mapping is necessary before finalizing the behavioural cohort and its links to questionnaires. The unchanged full-release calculations document the current data, not 42 verified independent recordings. Numeric and categorical coding, partial-session exclusion, participant-cluster-robust uncertainty, and equal-participant slopes led to the same qualitative conclusion, but none resolves this linkage issue. The intervals permit effects in either direction, so the result does not establish equivalence or the absence of respiratory influences on PPS.

Audio--Tactile responses were faster than tactile-only responses on average, but neither the raw SOA gradient nor subtraction of Baseline at the same nominal delay label isolates auditory proximity. Historical source checks found unequal trial-relative tactile-cue times for equal labels; the public tables cannot rerun those waveform checks. Temporal expectancy, response probability, trial counts, and control conditions matter when interpreting near-versus-far tactile benefits \citep{holmes2020sounds}. Baseline phase contrasts were uncertain after adjustment across the supplementary family; these data establish neither phase invariance nor an adjusted Baseline phase effect.

The participant-level sigmoid gate failed in every trial cell. The saturated five-SOA ceiling shows that the required trial-level variance explanation is unattainable for any deterministic curve of these SOAs, independently of nonlinear optimizer settings. Better in-sample sigmoid than linear RMSE therefore does not supply valid participant-boundary estimates. Four parameters fitted to five SOA means leave one residual degree of freedom and often meet parameter limits; group curves likewise cannot provide participant replicates. The disagreement between approximate registered distances and the stated trajectory changes some descriptive fit diagnostics but leaves the trial-level gate failure unchanged. Participant-boundary, sharpness, and related questionnaire hypotheses consequently remain unestimated under the planned gate.

The separate descriptive extension obtained four sigmoid approximations through a uniform signal-extraction and smoothing procedure. The original 320-pipeline winner yielded broad transitions whose 10--90\% spans all exceeded the measured range. Requiring a visibly S-shaped transition produced four sharper approximations, but this geometry was imposed: three widths reached their floor and one midpoint reached the sampling limit. In nested participant prediction, the constrained family slightly improved on equally processed lines but not on a constant. Unsmoothed curvature tests likewise did not establish nonlinearity. Neither result proves that the population response is linear or constant; together they leave a data-supported PPS sigmoid unresolved. The midpoint and slope comparisons are therefore assumption-dependent descriptions. Conditional bootstrap intervals omit method-selection uncertainty and are restricted by the imposed geometry. A future design with more positions spanning both plateaus, suitable expectancy controls, and an independently specified extraction/bandwidth rule would allow a stronger test of midpoint and sharpness differences. The later 40-record exclusion audit compared that logarithmic alternative under the same nested target. The visible sigmoid narrowly outperformed the constant by 0.105~ms, while the log-distance family did not improve prediction. This small descriptive difference does not establish an S-shaped mechanism. Nine prespecified-for-this-extension geometry specifications exposed substantial sensitivity under a permissive rule, and 30 repeated selections chose 27 distinct pipelines. These findings strengthen the need to separate numerical approximation from identified physiological parameters. They do not replicate the expectancy-controlled experiment of \citet{hobeika2019capturing}, calibrate method-selection-adjusted intervals or exhaust plausible shape definitions.

The individual-parameter extension achieved a computational description for all four cells of every participant record, including normalized common plateaus. It did not establish reliable individual boundaries. Shared plateaus slightly improved prediction and some recovery measures, but midpoint split-half reliability remained low, uncertainty was wide and the main questionnaire links changed with plateau assumptions. Every main refit-bootstrap interval included zero. The penalty can stabilize numerical fitting while shrinking genuine differences; unit-amplitude normalization can make negligible or uncertain changes look prominent. The raw-RT panels, direction flags and amplitude intervals therefore remain necessary companions to the normalized S shapes. Partial pooling and recovery studies provide a methodological rationale \citep{Ahn2011}, while the distinction between group effects and individual-difference reliability \citep{Hedge2018} explains why attractive curves alone cannot support correlations. The three recording-exclusion analyses now include the primary RT models and refit intervals for parameter contrasts and questionnaire links. The primary RT conclusion was unchanged; all 72 mean-contrast and all 216 main association intervals across cohorts and plateau models included zero. Midpoint-change split-half correlations were approximately zero, directly limiting the use of those changes in individual-difference correlations. These analyses expose sensitivity but do not recover the missing source mapping. These post hoc proxies cannot substitute for the failed preregistered boundary analysis.

The questionnaire extension adds a post-intervention experiential outcome alongside the behavioural findings. Paired spatial-frame codes tended to be higher under Viscereality, despite both condition medians being D. The owner confirmed that MAIA-2 preceded the experiment, supporting its use as a baseline trait correlate/moderator. Higher Trusting was associated with a larger signed Viscereality--Control code difference. Its positive association appeared in numerical-delta and direction-only correlations and in Condition $\times$ MAIA models with ordinary and HC3 uncertainty, with and without presentation order. These analyses reuse the same participants and do not constitute independent replications. Noticing was less consistent: HC3 interaction inference differed from the ordinary interaction and main rank analyses. The other dimensions did not show clear numerical-delta associations, which does not establish independence. Numerical differences assume comparable category spacing; the direction-only sensitivity avoids that spacing assumption. Without the administered pictogram key, higher codes cannot yet be assigned an allocentric or egocentric interpretation. Order adjustment does not remove all confounding or resolve carryover. These post hoc, unadjusted findings describe associations and trait moderation, not mediation or a PPS-boundary mechanism. MAIA-2 self-report is conceptually distinct from objective interoceptive accuracy \citep{Mehling2018,Garfinkel2015}. The small sample, uncertainty in reliability and multiple comparisons limit inference. Original questionnaire registration wording, administered anchors and language still require verification before submission.

Several features constrain physiological interpretation. Phase labels identify the instructed half-cycle rather than measured, trial-by-trial belt phase. Historical acquisition sources differ in RT definitions, available sounds, and recording availability; the released CSVs do not identify versions or support uniform waveform harmonisation. The hit-probability and time-on-task checks do not identify RTs on missed trials. The missing-hit-RT scenarios probe specified departures from cell medians and cannot establish the missingness mechanism. The recorded behavioral results should therefore not be read as a precise test of physiologically measured PPS boundaries.

Registration distinctions also require transparent reporting. The Bayesian mixed-model comparison was not fully specified, so the exploratory paired JZS factor addresses a different model and is sensitive to its normal-likelihood assumption. The sigmoid exclusion's scope and computational details required explicit interpretation. The registration disclosed pilot observations, and the released roster exceeds its stopping target; the collection history and ethics-statement details require author documentation before submission. Those items cannot be inferred from arbitrary public participant IDs.

The marimo publication companion makes the data-supported analysis inspectable and executable in the reader's browser. It exposes preprocessing counts, fixed analysis choices, model uncertainty, failed fit gates, unavailable provenance checks, and downloadable source and result hashes. Independent native checks support numerical agreement without substituting precomputed results for browser execution. This reproducibility supports assessment of the reported null detection while preserving the distinction between an executable analysis draft and a paper ready for submission.
